\chapter{Phase 2 --- Topographic Illumination Coupling}
\label{chap:phase2}

\begin{milestone}
\textbf{Status:} in progress.\\
\textbf{Acceptance:} with a local slope $s=\SI{6}{\degree}$ and
pole-facing aspect, the ChaSTE peak-surface-temperature bias of
$\sim\SI{-40}{\kelvin}$ identified in \cref{fig:chaste-val} is reduced
to $|{\rm bias}|\le\SI{10}{\kelvin}$.\\
\textbf{Deliverable:} \code{notebooks/phase2\_illumination/01\_slope\_horizon\_shadow.ipynb}
and \code{02\_chaste\_slope\_rerun.ipynb}.
\end{milestone}

\section{Goal of Phase 2}

Phase~1 proved that the vertical physics is correct.  Phase~2 closes
the last degree of freedom in the \emph{horizontal} boundary condition
by coupling each solver column to its local DEM context: slope,
aspect, horizon, and sky view factor.  This is the minimum machinery
required to (i)~reproduce the ChaSTE point measurement and (ii)~make
Phase~3's polar mapping meaningful.

\section{Pipeline data flow}

\begin{figure}[H]
\centering
\begin{tikzpicture}[node distance=0.9cm and 1.4cm,
    every node/.style={phaseBlock, minimum width=3.1cm, minimum height=1.0cm}]
  \node                                       (dem)   {DEM mosaic\\\tiny LOLA\,+\,Kaguya};
  \node[right=of dem]                         (slope) {Slope/aspect\\\tiny $s,\varphi_s$};
  \node[right=of slope]                       (hor)   {Horizon raycast\\\tiny $h(\phi)$};
  \node[below=of hor]                         (sky)   {Sky view factor\\\tiny $F_{\mathrm{sky}}$};
  \node[left=of sky]                          (shad)  {Shadow mask\\\tiny $m(t)$};
  \node[left=of shad]                         (ins)   {Corrected $Q(t)$\\\tiny \cref{eq:Qtotal}};
  \node[below=of ins, activeBlock, minimum width=3.1cm] (cn) {Phase-1 solver\\\small (reused)};
  \draw[arrowThick] (dem)   -- (slope);
  \draw[arrowThick] (slope) -- (hor);
  \draw[arrowThick] (hor)   -- (sky);
  \draw[arrowThick] (hor)   -| (shad);
  \draw[arrowThick] (slope.south) |- (shad.east);
  \draw[arrowThick] (sky)   -- (shad);
  \draw[arrowThick] (shad)  -- (ins);
  \draw[arrowThick] (ins)   -- (cn);
\end{tikzpicture}
\caption[Phase 2 data flow]{Phase 2 data flow.  The Phase-1 solver is
  reused unchanged; only its surface forcing is replaced by the
  topographically-corrected $Q(t)$.}
\label{fig:phase2-flow}
\end{figure}

\section{Digital elevation model}

The preferred DEM is the merged LOLA/Kaguya product
\citep{barker2016,araki2009} at \SI{5}{\meter\per pixel} polar
stereographic, cropped to a \SI{10}{\kilo\meter} square around each
target column.  Loading is delegated to \code{rasterio}:

\begin{lstlisting}[caption={DEM loader from \code{lunar/dem.py}.},label={lst:dem_load}]
import rasterio
from rasterio.warp import transform

def load_dem_window(path, centre_lonlat, half_size_m=5000.0, resolution=5.0):
    """Load a square window of the DEM centred on (lon, lat)."""
    with rasterio.open(path) as src:
        xs, ys = transform("EPSG:4326", src.crs,
                           [centre_lonlat[0]], [centre_lonlat[1]])
        row, col = src.index(xs[0], ys[0])
        half_px  = int(half_size_m / resolution)
        window   = rasterio.windows.Window(
            col - half_px, row - half_px, 2 * half_px, 2 * half_px)
        elev = src.read(1, window=window)
        transform_w = src.window_transform(window)
    return elev, transform_w
\end{lstlisting}

\section{Slope and aspect}

The local slope magnitude $s$ and aspect $\varphi_s$ are obtained from
the DEM by central-difference gradients:
\begin{equation}
  s = \arctan\!\sqrt{(\partial_x z)^2 + (\partial_y z)^2},
  \qquad
  \varphi_s = \mathrm{atan2}(-\partial_y z,\ \partial_x z).
  \label{eq:slope_aspect}
\end{equation}

\section{Horizon raycast}

For each DEM pixel, the horizon elevation $h(\phi)$ is the maximum
apparent elevation angle sampled along $N_\phi$ azimuth rays
($N_\phi = 360$ by default):
\begin{equation}
  h(\phi) \;=\; \max_{r \in (0, r_{\max}]}
  \arctan\!\left(\frac{z(\mathbf{x} + r\,\hat{\bm{n}}_\phi) - z(\mathbf{x})}{r}\right),
  \label{eq:horizon}
\end{equation}
where $\hat{\bm{n}}_\phi = (\cos\phi, \sin\phi)$.  The raycast is
embarrassingly parallel across pixels and is GPU-ready; a Numba
CPU implementation is sufficient for single-column work.

\begin{lstlisting}[caption={Horizon raycast kernel (sketch) from \code{lunar/horizon.py}.},label={lst:horizon}]
@njit(cache=True, parallel=True)
def horizon_profile(elev, i0, j0, n_phi=360, r_max=1000):
    """Return horizon elevation (rad) at pixel (i0, j0) for n_phi azimuths."""
    H = np.empty(n_phi)
    for k in prange(n_phi):
        phi = 2.0 * np.pi * k / n_phi
        dx, dy = np.cos(phi), np.sin(phi)
        h_max = -np.pi / 2.0
        for r in range(1, r_max + 1):
            ii = int(round(i0 + r * dy))
            jj = int(round(j0 + r * dx))
            if ii < 0 or ii >= elev.shape[0] or jj < 0 or jj >= elev.shape[1]:
                break
            dz = elev[ii, jj] - elev[i0, j0]
            h = np.arctan2(dz, float(r))           # pixel distance -> metres upstream
            if h > h_max:
                h_max = h
        H[k] = h_max
    return H
\end{lstlisting}

\section{Sky view factor}

The diffuse-sky multiplier on the infrared background and the scattered
solar multiplier on albedo-reflected photons share the same geometric
kernel: the fraction of the hemisphere not occluded by the surrounding
terrain.  Under the flat-sky Lambertian approximation
\citep{dozier1990},
\begin{equation}
  F_{\mathrm{sky}}
  \;=\;
  \frac{1}{2\pi}\int_{0}^{2\pi}
  \left[\,\cos^{2}\!\max(h(\phi),0)\,\right]\,d\phi.
  \label{eq:Fsky}
\end{equation}
For a flat surface $h(\phi)\equiv 0$ and $F_{\mathrm{sky}}=1$; for the
ChaSTE site with $s=\SI{6}{\degree}$ and the surrounding crater rim,
$F_{\mathrm{sky}}\approx 0.94$, a \SI{6}{\percent} reduction that
cools the night-side temperature by $\sim\SI{2}{\kelvin}$.

\section{Shadow mask}

Given the instantaneous solar zenith $\beta(t)$ and azimuth
$\alpha(t)$ from SPICE, a pixel is in shadow when
\begin{equation}
  \beta(t) > \pi/2 - h(\alpha(t)).
  \label{eq:shadow}
\end{equation}
This is checked once per solver step and multiplies the direct-beam
term in $Q(t)$ by a binary mask.

\section{Corrected surface forcing}

The full corrected top-of-regolith flux seen by the Phase-1 solver is
\begin{equation}
  Q(t) =
  m(t)\,\Rsolar \,\cos i(t)
  \,+\, F_{\mathrm{sky}}\,Q_{\mathrm{IR,sky}}
  \,+\, A\,\Rsolar\,\cos i(t)\,(1-F_{\mathrm{sky}}),
  \label{eq:Qtotal}
\end{equation}
with $\cos i$ given by \cref{eq:cos_incidence}, $m(t)$ by
\cref{eq:shadow}, and the third term a first-order treatment of
reflected-light enhancement from surrounding rims.  Inserting
\cref{eq:Qtotal} into \cref{eq:surface_bc} is a one-line change in the
solver; the rest of the Phase-1 kernel is untouched.

\begin{novelty}
The \emph{same} 1-D Crank--Nicolson kernel that passed Phase~1 accepts
the topographically-corrected forcing \cref{eq:Qtotal} without any
architectural change.  This preserves the Phase-1 acceptance tests as
a regression suite for Phase~2, and is only possible because the
solver is written around the abstract \code{InsolationDriver}
interface rather than a hard-coded flat-surface kernel.
\end{novelty}

\section{Validation plan}

\begin{enumerate}
  \item \textbf{Unit tests.}  \code{tests/test\_horizon.py} compares
        the raycast against analytic limits (pyramid, cone, hemisphere).
  \item \textbf{Flat-surface regression.}  With $s=0$ and a flat DEM,
        \cref{eq:Qtotal} must reduce to the Phase-1 flat forcing to
        machine precision.  This guarantees Phase~1 acceptance is not
        broken.
  \item \textbf{ChaSTE slope re-run.}  With $s=\SI{6}{\degree}$ and
        pole-facing aspect at \SI{69}{\degree}S, the peak-surface
        temperature bias must fall below \SI{10}{\kelvin}.
  \item \textbf{Apollo~15/17 re-run.}  With their near-zero local
        slopes, the Apollo RMSE values from \cref{tab:phase1-rmse}
        must shift by less than \SI{0.1}{\kelvin}.
\end{enumerate}

\section{Bridge to Phase 3}

With pixel-level topographic forcing in place, the solver can be swept
across a DEM mosaic to build a \emph{map} of equilibrium temperatures
and --- for the first time in the pipeline --- ice-stability depths.
\Cref{chap:phase3} extends the kernel with the ice-coupled
conductivity closure (\cref{eq:k_mix}) and Diviner validation.
